\section{Problem Framing}
\label{sec:problem}

\textbf{The challenge.} We must decide which transactions deserve attention when no labels exist, and explain
each decision, because the user acts on the reason, not on a score.

\textbf{Limits of existing solutions.} Block explorers list transactions without ranking them. A fixed rule, such
as ``flag anything above one bitcoin'', ignores scale: wallets, trading bots and cold storage behave very
differently, and the median amounts of our addresses differ by a factor of \V{eda.scale_ratio}. One threshold
over-flags busy addresses, misses changes on quiet ones, and cannot follow changing behaviour. Supervised models
trained on public labelled datasets learn from sparse labels and from aggregated or anonymised features that cannot
be computed for an arbitrary address (Akcora et al.\ 2020; Elmougy and Liu 2023), and their accuracy collapses when
behaviour changes (Pareja et al.\ 2020).

\textbf{Why machine learning.} Learning can model each address's normal behaviour, find structure without
labels, read a history as a sequence, and use labels where they exist (Hilal, Gadsden and Yawney 2022; Han et
al.\ 2022). Real transactions have no target label, so we combine unsupervised methods with regression and use
classification only on labelled data. Four research questions guide the work:
\begin{enumerate}[label=RQ\arabic*]
  \item Which stable method best finds transactions planted in an address's own history?
  \item How well do classifiers perform on labelled data over time and on unseen families?
  \item Does ``unusual'' coincide with the labelled datasets' own illicit class?
  \item Does a sequence model trained on other addresses generalise better than per-address statistics, and can a
    language model trained from scratch explain flags faithfully?
\end{enumerate}
